\documentclass[10pt,a4paper]{amsart}
\usepackage[margin=19mm]{geometry}
\usepackage{amsmath,amssymb,mathtools}
\usepackage[scr=rsfs,cal=euler]{mathalfa}
\usepackage{xfrac}
\usepackage[unicode,hidelinks,bookmarks=false]{hyperref}
\newcommand{\ie}{i.e.\ }
\newcommand{\cf}{cf.\ }
\newcommand{\resp}{respectively\ }
\newcommand{\Subsection}{\subsection}
\providecommand{\nobreakdash}{\nobreak\hbox{-}}
\setlength{\emergencystretch}{2em}
\multlinegap0pt
\usepackage{amssymb}
\usepackage{xcolor}
\usepackage{float}
\newtheorem{lemma}{Lemma}
\newcommand{\BP}{\operatorname{BP}}
\title{The variation of zeros of the Miller basis}
\author{Liubomir Chiriac, Andrei Jorza}
\date{Source: arXiv:2605.09731v1 (CC BY 4.0)}
\begin{document}
\maketitle

\noindent\emph{Source and licence.} The variation of zeros of the Miller basis. Version arXiv:2605.09731v1 (CC BY 4.0), distributed by arXiv under the Creative Commons Attribution 4.0 International licence, http://creativecommons.org/licenses/by/4.0/, as recorded in the arXiv licence metadata for this exact version and shown on the abstract page https://arxiv.org/abs/2605.09731v1. Full licence notice: \texttt{licenses/arxiv-2605.09731-CC-BY-4.0.txt}.\par\smallskip

\noindent\textit{Editorial scope.} Counted unit: the article's lemma l:BP (source0.tex lines 1330--1345). The article's complete original proof of it is delivered with it (source0.tex lines 1346--1509, one \texttt{proof} environment of 164 lines). No mathematics is written, repaired or re-derived: the statement and the proof are the article's own lines, in the article's own notation, and no retained line is substituted or re-flowed.  Retained uncounted as the source-local context the proof needs: lines 1305--1318.  One declared adaptation: exactly 4 ancillary strings inside the retained spans are omitted (image-inclusion commands and, where the source repeats a label, the repeated label definitions) and no other line differs from the source; the figure environments, with their placement, captions and labels, are kept, and the package records the omitted strings in fidelity receipts bound to the affected bodies.  No retained line contains a citation, so the wrapper carries no bibliography and no bibliography entry.  The retained lines contain the article's own diagram code, which the wrapper typesets with the standard tikz packages; no image file is delivered and the package declares no asset dependency.  Editorial formatting only, in addition: an AMS \texttt{amsart} preamble that restates the article's own declarations and macros used by the retained lines, namely the environments \texttt{lemma} on separate counters, together with the standard packages those lines need.  The retained environments print in this document as Lemma 1 in order of appearance, whereas the article prints the same items with the numbering of its own full document.  The complete native source of the article as distributed in the arXiv e-print for this exact version is bundled unchanged as \texttt{sources/200249/source0.tex}.
\begin{lemma}\label{l:BPdiff}
Given $z< x$, the function $h(y)=\BP(x,y)+\BP(y,z)$ on
the interval $(z,x)$ has the following monotonicity
\begin{center}
\begin{tabular}{llll}
  $(x,y)$&$(y,z)$&Increasing&Descreasing\\
  \hline
  good ($y\geq cx$)&good ($z\geq cy$)&$y< \frac{x+z}{2}$&$y> \frac{x+z}{2}$\\
  good ($y\geq cx$)&bad ($z< cy$)&$y< \frac{x}{2-c}$&$y> \frac{x}{2-c}$\\
  bad ($y< cx$)&good ($z\geq cy$)&always&\\
  bad ($y< cx$)&bad ($z< cy$)&always&
\end{tabular}
\end{center}
\end{lemma}

\begin{lemma}\label{l:BP}
If $D-j=i_s<i_{s-1}<\ldots<i_0=D-i$ then
\[\prod_{r=0}^{s-1}c_{i_r,-i_{r+1}}\leq \max(T_1,T_2,T_2',T_3,T_3',T_4),\]
where
\begin{align*}
  T_1&=(1728-e^{2\pi})^{j-i}s^{j-i}\frac{(D-i)^{D-i}}{(D-j)^{D-j}(j-i)^{j-i}}\\
  T_2&=2981^{D-i}1.25^{D-j}\\
  T_2'&=(2393s)^{D-i}1.45^{j-i}\\
  T_3&=(2386s)^{D-i}2924^{-(D-j)}\\
  T_3'&=5849^{D-i}2924^{-(D-j)}\\
  T_4&=1728^{D-i}e^{-2\pi (D-j)}.
\end{align*}
When $D-j>c(D-i)$, it suffices to consider only $T_1$. In
particular, if $\mathcal{A}=5849$ and $\mathcal{B} = 1.5$ then
\[\prod_{r=0}^{s-1}c_{i_r,-i_{r+1}}\leq \max(T_1,(\mathcal{A}s)^{D-i}\mathcal{B}^{D-j}).\]
\end{lemma}

\begin{proof}
We will bound
\[\sum_{r=0}^{s-1} \log c_{i_r,-i_{r+1}}\leq
\sum_{r=0}^{s-1}\BP(i_r,i_{r+1}).\]
To estimate an upper bound, we will determine the maximum of the continuous function
\[\BP(x_0,x_1,\ldots,x_n)=\sum_{r=0}^{s-1}\BP(x_r,x_{r+1}),\]
in the compact set
\[\{(x_0,\ldots, x_n)\mid A=x_n\leq x_{n-1}\leq\ldots \leq x_0=B\}.\]
Let $(x_0,\ldots, x_n)$ be the point
where the maximum is attained.

Since $\BP(x,x)=0$, we may instead assume that the maximum is
attained at a point $(y_0,\ldots, y_m)$ such that
$A=y_m<y_{m-1}<\ldots<y_0=B$, where $m\leq n$, by eliminating
duplicates.

We will represent the sequence $A=y_m<y_{m-1}<\ldots <y_0=B$ as
a ladder

% sage: def ladder(xs,labels,eps=0.1,textx=0.5, texty=0.5, fontsize='large', figsi
% ....: ze=5, dpi=300):
% ....:     X = max(xs)*(1+eps)
% ....:     G=line([(0,0),(X,X)],aspect_ratio=1, ticks=[[],[]], figsize=figsize, d
% ....: pi=dpi)
% ....:     for i in range(len(xs)-1):
% ....:         G= G+line([(xs[i],xs[i]),(xs[i+1],xs[i])])
% ....:         G= G+line([(xs[i+1],xs[i]),(xs[i+1],xs[i+1])])
% ....:     for i in range(len(xs)):
% ....:         G= G+text("{}".format(labels[i]), (xs[i]-textx,xs[i]+texty), fonts
% ....: ize=fontsize)
% ....:     return G

% sage: ladder([1,1.5,2.3,4,5,7,10],['$A$','','$\\ldots$','$\\ldots$', '$y_2$','$y
% ....: _1$','$B$'], fontsize='x-large', figsize=5)+text("$(y_k,y_{k+1})$", (4.5,1.8), fontsi
% ....: ze='x-large')

\begin{center}

\end{center}

To describe $(y_0,\ldots, y_m)$ at the maximum, we will
repeatedly use Lemma \ref{l:BPdiff}. We denote $\mathcal{L}$ the line $y=cx$ and $\mathcal{L}'$ the line $y=\frac{1}{2-c}x$.

Suppose that $P_k=(y_k,y_{k+1})$ is bad. First, note that $P_{k-1}=(y_{k-1},y_k)$
can't also be bad, since increasing
$y_k$ would increase the function. Thus 
$P_{k-1}$ is good and located on the line $\mathcal{L}'$, or else increasing or decreasing
$y_k$ would increase the function.

Suppose $P_k=(y_k,y_{k+1})$ is good, but not on the line
$\mathcal{L}$. If $P_{k-1}$ is bad, we could increase $y_k$ to increase
the value of the function, which means $P_{k-1}$ is also
good. If $P_{k-1}$ is on the line $\mathcal{L}$, then
$y_{k-1}+y_{k+1}>\frac{1}{c}y_k>2y_k$, so we may increase $y_k$
keeping both $P_k$ and $P_{k-1}$ good to increase the value of
the function. If $P_{k-1}$ is not on the line $\mathcal{L}$, it must be
that $2y_k=y_{k-1}+y_{k+1}$. We conclude that every good point not on the line $\mathcal{L}$ must be preceeded by evenly spaced out good points, also not on the line $\mathcal{L}$.

Finally, if $P_k$ is on the line $\mathcal{L}$, by the same
argument, the previous point $P_{k-1}$ must be good with
$y_{k-1}-y_k\geq y_k-y_{k+1}$.

Therefore, there are four possible configurations:
\begin{enumerate}
\item All points $P_0,\ldots, P_{m-1}$ are good, not on
$\mathcal{L}$, with $y_0,\ldots, y_m$ evenly spaced out.
\item The last point $P_{m-1}$ is bad, $P_{m-2}$ is on
$\mathcal{L}'$, and $y_0,\ldots, y_{m-1}$ are evenly spaced out.
\item The points $P_0,\ldots, P_{q-1}$ are good, not on
$\mathcal{L}$, with $y_0,\ldots, y_q$ evenly spaced out, and
$P_q,\ldots, P_{m-1}$ are all on $\mathcal{L}$.
\item The last, degenerate, case is $m=1$, $y_0=B$, $y_1=A$.
\end{enumerate}

% sage: ladder([3,4,5,6,7,8,9],['$A$','','$\\ldots$','$\\ldots$', '$y_2$','$y_1$',
% ....: '$B$'], fontsize='x-large', figsize=5)+line([(0,0),(10,5)],color="red")+te
% ....: xt("$\\mathcal{L}$", (6,2.3), fontsize='x-large')

% sage: ladder([1,4,5,6,7,8,9],['$A$','','$\\ldots$','$\\ldots$', '$y_2$','$y_1$',
% ....: '$B$'], fontsize='x-large', figsize=5)+line([(0,0),(10,5)],color="red")+li
% ....: ne([(0,0),(10,8)],color="purple")+text("$\\mathcal{L}$", (6,2.3), fontsize
% ....: ='x-large')+text("$\\mathcal{L}'$", (9,6), fontsize='x-large')

% sage: ladder([0.7,0.7*2,0.7*4,9-1.55*3,9-1.55*2,9-1.55,9],['$A$','','$y_q$','$\\
% ....: ldots$', '$y_2$','$y_1$','$B$'], fontsize='x-large', figsize=5)+line([(0,0
% ....: ),(10,5)],color="red")+text("$\\mathcal{L}$", (6,2.3), fontsize='x-large')


\begin{figure}[htp]
\centering
\hfill
\hfill

\end{figure}

In each of these cases, we can determine $y_0,\ldots, y_m$
exactly. For convenience, we denote $[x,y]=x\log x-y\log y-(x-y)\log(x-y)\leq x\log 2$.
\begin{enumerate}
\item For $1\leq u\leq m$, $y_i = B-(B-A)i/m$. In this case, the
value of the function is
\begin{align*}
\BP(y_0,\ldots,y_m)&=\sum_{i=0}^{m-1}\left((y_i-y_{i+1})\log(1728(1-c))+y_i\log
                     y_i-y_{i+1}\log y_{i+1}-\frac{B-A}{m}\log
                     \frac{B-A}{m}\right)\\
  &=(B-A)(\log m+\log(1728(1-c)))+B\log B-A\log A-(B-A)\log (B-A)\\
  &=(B-A)(\log m+\log(1728(1-c)))+[B,A],
\end{align*}
which is maximized when $m=n$.
\item The remaining three cases can only occur if $A< \mathcal{B}$. We have $y_m=A$, $y_{m-1}=A+a$, $y_{m-2}=A+a+b$, \ldots,
$y_0=B=A+a+(m-1)b$ with $y_{m-1}=\frac{1}{2-c}y_{m-2}$. This
gives $b=(1-c)(A+a)$ and
\[a=\frac{B+A}{1+(m-1)(1-c)}-A.\]
In this case, denoting $N=1+(m-1)(1-c)$ the
value of the function is
\begin{align*}
g(A+a,A) &+\BP(y_0,\ldots,y_{m-1})\\ &=(A+a)\log 1728-2\pi
                                  A+(B-A-a)(\log(m-1)+\log(1728(1-c)))\\
                                &\ \ +B\log B-(A+a)\log (A+a)-(B-A-a)\log (B-A-a)\\
                                &=B(\log(1728)+\log(N-1))-(A+a)\log
                   (N-1)+[B,A+a]\\
  &\leq B(\log 1728+\log 2)+(B-\frac{A+B}{N})\log(N-1).
\end{align*}
When $m=2$, $N=2-c$ and the quantity above is
\[\leq B\left(\log 1728+\log 2+\left(1-\frac{1}{2-c}\right)\log(1-c)\right)-\frac{\log(1-c)}{2-c}A\leq 8B+0.22A.\]
When $m\geq 2$, $N=1+(1-c)(m-1)>2$ and the quantity above is
maximized when $m=n$ giving
\[\leq B\left(\log 1728+\log
2+\left(1-\frac{2}{N}\right)\log(N-1)\right)+(B-A)\frac{\log(N-1)}{N-1}\]
\[<B\left(\log 1728+\log
2+\log(1-c)+\log(n)\right)+e^{-1}(B-A)<(7.78+\log n)B+e^{-1}(B-A).\]



\item For $0\leq i\leq m-q$ we have $y_{m-i}=c^{-i}A$, and $y_j=B-(B-c^{q-m}A)j/(m-q)$ for $0\leq j\leq q$. In this case, the
value of the function is
\[A \left(\frac{c^{q-m}-1}{c^{-1}-1}\right)\left(c^{-1}\log
1728-2\pi\right)+(B-c^{q-m}A)(\log(q)+\log(1728(1-c)))+\left[B,c^{q-m}A\right]\]
\[\leq B\log 2 + A \left(\frac{c^{q-m}-1}{c^{-1}-1}\right)\left(c^{-1}\log
1728-2\pi\right)+(B-c^{q-m}A)(\log(q)+\log(1728(1-c)))\]
\[= B\log 2 - A \left(\frac{\log
  1728-2\pi c}{1-c}\right)+B(\log(q)+\log(1728(1-c)))+\]
\[+c^{q-m}A\left(\frac{(\log
1728-2\pi) c}{1-c}-(\log(q)+\log(1-c))\right)\]
Note that the last term is negative if $q\geq 3$. Keeping in
mind that $Ac^{q-m}\leq B$ we get that this is
\[\leq B\log 2 - A \left(\frac{\log
1728-2\pi c}{1-c}\right)+B(\log(q)+\log(1728(1-c)))\]
\[\leq B\left(\log 2 +\log(1728(1-c))+\log n\right)- A \left(\frac{\log
1728-2\pi c}{1-c}\right)\leq (7.78+\log n)B-7.98A\]
if $q\geq 3$ and 
\[\leq B\left(\log 2+\frac{\log
1728-2\pi c}{1-c}\right) - A \left(\frac{\log
1728-2\pi c}{1-c}\right)\leq 8.68B -7.98A\]
if $q\leq 2$. 


We remark that this case can only occur if $A\leq cB$.
\item In the degenerate case when $y_0=B$ and $y_1=A$ the value
of the function is
\[B\log 1728-2\pi A.\]
\end{enumerate}


\end{proof}

\end{document}
